\documentclass[12pt, letterpaper]{article}
\date{\today}
\usepackage[margin=1in]{geometry}
\usepackage{amsmath}
\usepackage{hyperref}
\usepackage{cancel}
\usepackage{amssymb}
\usepackage{fancyhdr}
\usepackage{pgfplots}
\usepackage{booktabs}
\usepackage{pifont}
\usepackage{amsthm,latexsym,amsfonts,graphicx,epsfig,comment}
\pgfplotsset{compat=1.16}
\usepackage{xcolor}
\usepackage{tikz}
\usetikzlibrary{shapes.geometric}
\usetikzlibrary{arrows.meta,arrows}
\newcommand{\Z}{\mathbb{Z}}
\newcommand{\N}{\mathbb{N}}
\newcommand{\R}{\mathbb{R}}
\newcommand{\Q}{\mathbb{Q}}
\newcommand{\C}{\mathbb{C}}
\newcommand{\F}{\mathbb{F}}

\newcommand{\Po}{\mathcal{P}}
\newcommand{\Pro}{\mathbb{P}}
\author{Alex Valentino}
\title{Numerical Analysis homework}
\pagestyle{fancy}
\renewcommand{\headrulewidth}{0pt}
\renewcommand{\footrulewidth}{0pt}
\fancyhf{}
\rhead{
	Homework \\
	Numerical Analysis
}
\lhead{
	Alex Valentino\\
	Collaborators: Irshad Guliyev, Jay Patwardhan
}
\begin{document}
\begin{enumerate}
	\item 
	\begin{enumerate}
		\item 38.625 can be represented as $1*1.00110101*2^5$
		\item -10.25 can be represented as $-1*1.01001*2^3$
		\item 5.6 can be represented as $1*1.01\overline{1001}*2^2$
	\end{enumerate}
	Note that all of the answer above were arrived at by manually dividing the given answers by the largest possible power of 2, converting that
	into binary, then shifting until appropriate.  
	\item Consider a system which can stores numbers in the following way: $\sigma * a * 2^e$ where 
	$\sigma \in \{-1,1\}, 1 \leq a \leq 2, -1022 \leq e \leq 1023$ and $a$ has 10 bits of precision (i.e. 
	$a = 1.a_1a_2a_3\ldots a_{10}$). 
	\begin{enumerate}
		\item 5.6 round up: $1 * 1.0110011010 *2^2$, this is because the expression with more bits of the mantissa is 1.01100110011001, thus 
		to round up we just add the machine epsilon to bump it up. 
		\item 5.6 round to nearest: $1 * 1.0110011010 *2^2$ Note that to round down we would have a mantissa of $1.0110011001$, and since 
		as shown above with the expanded mantissa that the first decimal in the expanded mantissa not found within the original mantissa is $1$, and there are additional bits after it implies that the rounded number is infact closer to the true value.
		\item -5.6 using round toward 0: $-1 * 1.0110011001 *2^2$.  Here I noted that rounding towards zero would in fact mean rounding up this negative number, or effectively rounding down the mantissa.  This is observed
		\item -5.6 using round down $-1 * 1.0110011010 *2^2$.  Since we're rounding down a negative number that means that the mantissa of the number is increasing.  Therefore we use the rounded up mantissa we've seen before.  
	\end{enumerate}
	\item Consider a very limited floating-point system that expresses numbers in the form $a*2^e$.  
	In this system: the significand $a$ follows the form $1.a_1a_2a_3$ where each $a_i \in \{0,1\}$ and
	$e \in \{-1,0,1\}$.  
	\begin{enumerate}
		\item What is the largest number that can be represented in this system?  That would have to be $3.75$ or 
		$11.11$ in binary.
		\item What is the smallest positive number that can be represented?  That would have to be $0.5$ or 
		$0.1000$ in binary
		\item What is the machine epsilon?  That is, the smallest $\epsilon > 0$ such that $1+\epsilon$ is 
		representable.  The $\epsilon$ would have to be $0.125$ or $0.001$ in binary.  This is the case because 
		after the assumed $1$ in the decimal, there are three decimals of precision.  Therefore one can choose the last 
		digit to be 1, yielding $1.001$
	\end{enumerate}
	\item In the Simpsons, Homer dreamed up the equation $1782^{12} + 1841^{12} = 1922^{12}$.  Did Homer disprover 
	Fermat's last theorem?
	\begin{enumerate}
		\item Compute $\sqrt[12]{1782^{12} + 1841^{12}}$ in matlab with format short.  What does matlab report, and what 
		happens if you put in format long?
		\begin{itemize}
			\item The short formatted answer is 1922
			\item The long formatted answer is 1921.999999955876
		\end{itemize}
		\item Determine the absolute and relative error of $1782^{12} + 1841^{12} \approx 1922^{12}$
			\begin{itemize}
				\item Absolute error: $-7.0021199503409775*10^{29}$, obtained via finding the difference between $1782^{12} + 1841^{12}$ and
				$1922^{12}$
				\item Relative error: $-2.755427056723427*10^{-10}$.  I simply divided the value given above by $1922^12$ (In the definitions you guys give there isn't absolute value signs, they can be applied if that's how you mean to define them)
			\end{itemize}
		\item prove that the equality is false.  We note that mod 2 that $x^2 \equiv x \mod{2}$ since $0^2 \equiv 0 \mod{2}$ and $1^2 \equiv 1 \mod{2}$.  Therefore $1782^{12} + 1841^{12} \equiv 1782 + 1841 \equiv 0 + 1 = 1 \mod{2}$.
		Additionally, $1922 \equiv 0 \mod{2}$.  Since the parity of both sides of the equality is unequal, then 
		the equality for the integers is false.  
	\end{enumerate}
	\item Let $x_A = 2.63$  If it was rounded to that amount, then $2.625 \leq x_T \leq 2.635$.  Therefore 
	$|x_T - x_A| \leq 0.005$.
	\begin{itemize}
		\item $\sin(x)$: For the absolute error $E$ we have that $$|E| = |\sin(x_A) - \sin(x_T)| \approx |\cos(2.63)|*|x_T - x_A| \leq  0.005*|\cos(2.63)| = 0.005*0.871 = 0.004355$$
		Note that the first approximation is gotten from the slides, and employs the mean value theorem $f(b) - f(a) = f'(c)(b-a)$, where
		$c \in (a,b)$, assuming the differentiability of $f$ on the previously stated interval and continuity of the smallest closed set 
		containing that interval.
		For the relative error we have to bound $\sin(x_t)$ from below.  Note that since $\frac{\pi}{2} \leq x_T \leq \pi$ that $\sin$ is decreasing on the interval which contains $x_T$.  Therefore $\sin(2.635) \leq \sin(x_T)$.  Thus 
		we have that $Rel(\sin(x_T)) = \frac{|E|}{|\sin(x_T)|} \leq \frac{0.004355}{0.4825} = 0.009025$.
		Therefore $-0.004355 \leq E \leq 0.004355$ and $-0.009025 \leq Rel(\sin(x_A)) \leq 0.009025$
		\item $\ln(x)$: For the absolute error $E$ we have that 
		$$
		|E| = |\ln(x_A) - \ln(x_T)| \approx |\frac{1}{x_A}||x_T - x_A| = \frac{1}{2.63}|x_T - x_A| \leq 0.001901
		$$
		since $\frac{d}{dx}\ln(x) = \frac{1}{x}$ by the mean value theorem.
		Since $\ln(x)$ is an increasing function we can lower-bound it by evaluating it on the lower bound for $x_T$, 
		$2.625$, $\ln(2.625) = 0.96508 \leq \ln(x_T)$.  Therefore the relative error is:
		$Rel(\ln(x_A)) = \frac{|E|}{|ln(x_T)|} \leq \frac{0.001901}{0.96508} = 0.0019697$.
		This gives us the following bounds:
		$-0.001901 \leq E \leq 0.001901$\\
		$-0.0019697 \leq Rel(\ln(x_A)) \leq 0.0019697$
	\end{itemize}
\end{enumerate}
\end{document}
